\documentclass[10pt]{article}

% ---------- Document Identity (update these only) ----------
\newcommand{\DocStage}{}
\newcommand{\DocTitle}{Your Road to Tier One SOF}
\newcommand{\ProgramName}{182-Week Civilian-to-Selection Readiness Pipeline}
\newcommand{\ProgramSubtitle}{}

% ---------- Page ----------
\usepackage[legalpaper,left=0.5in,right=0.5in,top=0.6in,bottom=0.45in]{geometry}

% ---------- Typography ----------
\usepackage[T1]{fontenc}
\usepackage[utf8]{inputenc}
\usepackage{libertinus}
\usepackage{microtype}
\usepackage{parskip}
\usepackage{ragged2e}
\usepackage{textcomp}
\AtBeginDocument{\color{black}}
\AtBeginDocument{\pagecolor{white}}
\AtBeginDocument{\justifying}

% ---------- Landscape pages ----------
\usepackage{pdflscape}

% ---------- Color ----------
\usepackage[table]{xcolor}
\definecolor{StageBlue}{HTML}{1F4E79}
\definecolor{StageBlueDark}{HTML}{163A5A}
\definecolor{LightGray}{HTML}{F2F4F7}
\definecolor{MidGray}{HTML}{D9DEE5}
\definecolor{RuleGray}{HTML}{C7CED8}
\definecolor{HeaderInk}{HTML}{000000}
\definecolor{Ink}{HTML}{000000}

% ---------- Utility ----------
\usepackage{enumitem}
\sloppy
\setlength{\emergencystretch}{2em}

% ---------- Boxes / UI ----------
\usepackage[most]{tcolorbox}

\tcbset{
  charterTitle/.style={
    enhanced,
    colback=StageBlue!4,
    colframe=StageBlue,
    boxrule=1.25pt,
    arc=3pt,
    left=16pt,right=16pt,top=14pt,bottom=14pt,
    borderline north={3.2pt}{0pt}{StageBlue},
    borderline west={4pt}{0pt}{StageBlue},
    borderline south={1.25pt}{0pt}{StageBlue},
    drop shadow={black!25!white},
  },
  charterRule/.style={
    enhanced,
    colback=LightGray,
    colframe=RuleGray,
    boxrule=0.85pt,
    arc=3pt,
    left=12pt,right=12pt,top=10pt,bottom=10pt,
    borderline west={3pt}{0pt}{StageBlue},
  },
  charterBadge/.style={
    enhanced,
    colback=StageBlue,
    coltext=white,
    colframe=StageBlueDark,
    boxrule=0.6pt,
    arc=2pt,
    left=6pt,right=6pt,top=3pt,bottom=3pt,
  }
}
\newtcbox{\Badge}{nobeforeafter,charterBadge}

% ---------- Lists ----------
\setlist[itemize]{leftmargin=1.5em, itemsep=0.25em, topsep=0.25em}
\setlist[enumerate]{leftmargin=1.8em, itemsep=0.25em, topsep=0.25em}

% ---------- TOC ----------
\usepackage{tocloft}
\setcounter{tocdepth}{3}
\setlength{\cftbeforesecskip}{3pt}
\setlength{\cftbeforesubsecskip}{2pt}
\setlength{\cftbeforesubsubsecskip}{0pt}
\renewcommand{\cftsecfont}{\bfseries}
\renewcommand{\cftsecpagefont}{\bfseries}
\renewcommand{\cftsubsecfont}{\normalfont}
\renewcommand{\cftsubsecpagefont}{\normalfont}
\renewcommand{\cftsubsubsecfont}{\normalfont}
\renewcommand{\cftsubsubsecpagefont}{\normalfont}
\setlength{\cftsecindent}{0pt}
\setlength{\cftsubsecindent}{1.25em}
\setlength{\cftsubsubsecindent}{2.8em}
\setlength{\cftsecnumwidth}{2.1em}
\setlength{\cftsubsecnumwidth}{2.8em}
\setlength{\cftsubsubsecnumwidth}{3.6em}

% Remove the built-in TOC heading so only the custom heading is shown
\makeatletter
\renewcommand{\tableofcontents}{\@starttoc{toc}}
\makeatother

% ---------- Header/Footer: page number only ----------
\usepackage{fancyhdr}
\pagestyle{fancy}
\fancyhf{}
\renewcommand{\headrulewidth}{0pt}
\renewcommand{\footrulewidth}{0pt}
\fancyfoot[C]{\thepage}

% ---------- Section formatting ----------
\usepackage{titlesec}

\titleformat{\section}
  {\Large\bfseries\color{Ink}}
  {\thesection}{0.6em}{}
  [\vspace{0.25em}\textcolor{StageBlue}{\rule{\linewidth}{0.8pt}}]

\titleformat{\subsection}
  {\large\bfseries\color{Ink}}
  {\thesubsection}{0.6em}{}

\titleformat{\subsubsection}
  {\normalsize\bfseries\color{Ink}}
  {\thesubsubsection}{0.6em}{}

\titlespacing*{\section}{0pt}{0.95em}{0.35em}
\titlespacing*{\subsection}{0pt}{0.65em}{0.25em}
\titlespacing*{\subsubsection}{0pt}{0.55em}{0.2em}

% ---------- Table engine ----------
\usepackage{tabularray}
\UseTblrLibrary{booktabs}

% ---------- tabularray: GLOBAL table treatment ----------
\SetTblrInner{
  cells = {fg=black, bg=white, valign=t},
}

% ---------- Header helpers ----------
\newcommand{\Hdr}[1]{\shortstack[c]{\strut #1\strut}}

% ---------- Lines ----------
\newcommand{\TitleLine}{\textcolor{StageBlue}{\rule{0.98\linewidth}{0.95pt}}}
\newcommand{\SubTitleLine}{\textcolor{RuleGray}{\rule{0.98\linewidth}{0.65pt}}}

% ---------- Continued on next page ----------
\newcommand{\ContinuedOnNextPageInline}{%
  \par
  \vfill
  \noindent\textcolor{RuleGray}{\rule{\linewidth}{1pt}}\vspace{-2pt}
  \noindent\hfill\textit{Continued on next page}\par\vspace{-10pt}
  \noindent\textcolor{RuleGray}{\rule{\linewidth}{1pt}}
  \clearpage
}

% ---------- PDF hyperlinks ----------
\usepackage[hidelinks]{hyperref}
\hypersetup{
  colorlinks=false,
  pdfborder={0 0 0},
  linktoc=all,
  pdftitle={\DocTitle},
  pdfauthor={\ProgramName}
}

% =========================================================
\begin{document}

\section*{7.14 — Environmental Apparel, Cold / Hot / Wet / Inclement Weather}
\phantomsection
\addcontentsline{toc}{section}{7.14 — Environmental Apparel, Cold / Hot / Wet / Inclement Weather}

\subsection*{7.14.1 Purpose \& Scope}
\phantomsection
\addcontentsline{toc}{subsection}{7.14.1 Purpose \& Scope}

This category covers:

\begin{itemize}
\item Cold-weather layering systems (base, mid, shell) that support safe outdoor execution in cold and wind.
\item Wet-weather and wind management systems (rain/wind shells, seam posture, hood/cuff sealing) that prevent wet-to-cold failure.
\item Hot-weather heat-management apparel and sun protection posture (breathable fabrics, sun coverage posture).
\item Hands/feet/head interface items that materially change exposure safety (gloves/mitts, neck gaiter, hat), without duplicating footwear systems.
\item Packable weather contingencies that preserve training continuity (storm kit posture), with staging cross-reference to Stage 7.10 and carry integration cross-reference to Stage 7.07.
\end{itemize}

This overview crosswalks Stage 6 timing to the environmental-apparel conditions required for safe outdoor execution and admissible evidence.

This category does not cover, except by cross-reference:

\begin{itemize}
\item Footwear and foot care: Stage 7.11
\item General daily apparel: Stage 7.15
\item Protective / visibility equipment: Stage 7.19
\item Logistics / staging systems: Stage 7.10
\item Hydration containers and prep systems: Stage 7.13
\end{itemize}

This overview does not provide programming, test protocols, calendars, medical guidance, or protective-visibility prescriptions that belong in 7.19.

\subsection*{7.14.2 Governance And Safety Boundaries}
\phantomsection
\addcontentsline{toc}{subsection}{7.14.2 Governance And Safety Boundaries}

\begin{itemize}
\item Stage 0 boundary alignment: stop rules, safety override doctrine, conflict-resolution hierarchy, and evidence-admissibility controls apply to all 7.14 selections and substitutions.
\item Environmental apparel is a safety control and a validity control: unsafe exposure can terminate sessions/tests and contaminate evidence.
\item Stop posture (non-medical and operational): substitute indoor modality or reschedule to the next valid window under hard-stop environment conditions. Do not push through.
\item Hard-stop environment conditions (examples):
  \begin{itemize}
  \item lightning or storm risk,
  \item ice glaze or otherwise unsafe footing,
  \item whiteout or near-zero visibility,
  \item official air-quality warning,
  \item extreme windchill concern,
  \item heat index concern.
  \end{itemize}
\item No novelty near gates:
  \begin{itemize}
  \item do not introduce new cold/wet/heat systems inside Deload+Test windows,
  \item log apparel system changes that materially affect comfort, pace, tolerance, or safety.
  \end{itemize}
\item Visibility and traffic safety: cross-reference Stage 7.19 only; do not duplicate Stage 7.19 content.
\end{itemize}

\subsection*{7.14.3 Tier Doctrine}
\phantomsection
\addcontentsline{toc}{subsection}{7.14.3 Tier Doctrine}

\begin{itemize}
\item Price band convention: \$ = minimal household-cost item; \$\$ = modest purpose-built item; \$\$\$ = expanded-capability or redundancy item.
\item N/A rule: use N/A only when a field is genuinely not applicable and omission does not obscure safety, maintenance, or phase-timing logic.
\item Substitution posture: substitutions are acceptable only when they preserve exposure safety, visibility of stop conditions, and comparability of outdoor execution; substitutes that increase exposure risk or invalidate weather comparability break intent.
\item Tier 0: household-only and opportunistic coverage; not sufficient for sustained outdoor execution in winter conditions; treated as a temporary constraint state.
\item Tier 1: minimum viable cold/wet/heat control enabling safe outdoor execution when required by Stage 6 timing, plus the ability to substitute indoors without creating unsafe exposure.
\item Tier 2: expanded modularity and packable contingencies that reduce canceled sessions and reduce wet-to-cold failure; includes traction aids and storm kit.
\item Tier 3: advanced / overbuilt redundancy and specialized severe-weather protection; travel / vehicle-staging posture to reduce variance near high-consequence phases.
\end{itemize}

\subsection*{7.14.4 Selection Criteria \& Fit Checks}
\phantomsection
\addcontentsline{toc}{subsection}{7.14.4 Selection Criteria \& Fit Checks}

Fit and safety checklist (operational):

\begin{itemize}
\item Layering compatibility:
  \begin{itemize}
  \item base/mid/shell must fit without compressing insulation or restricting breathing.
  \end{itemize}
\item Moisture management:
  \begin{itemize}
  \item fast drying, venting options, and sweat handling are prioritized to prevent wet-to-cold failure.
  \end{itemize}
\item Wind and wet performance posture:
  \begin{itemize}
  \item hood coverage, cuff sealing, hem sealing posture; seam quality posture; water repellency care posture.
  \end{itemize}
\item Dexterity and grip:
  \begin{itemize}
  \item glove systems must allow safe handling of gear (buckles, zippers) without exposing hands to frost risk.
  \end{itemize}
\item Thermal safety posture:
  \begin{itemize}
  \item avoid overheating under exertion; ability to vent is required.
  \end{itemize}
\item Trial-first rule:
  \begin{itemize}
  \item validate on short exposures before long sessions or any protected window; do not first-use a new system in Deload+Test weeks.
  \end{itemize}
\end{itemize}

\subsection*{7.14.5 Setup, Safety, and Anchoring}
\phantomsection
\addcontentsline{toc}{subsection}{7.14.5 Setup, Safety, and Anchoring}

Setup doctrine (no programming):

\begin{itemize}
\item Pre-exposure checks:
  \begin{itemize}
  \item classify the session as cold/wet/heat risk class; if hard-stop conditions exist, substitute indoors.
  \item visibility posture is governed by Stage 7.19 (cross-reference only).
  \end{itemize}
\item Layering setup:
  \begin{itemize}
  \item manage sweat: vent early; avoid saturating base layer.
  \item wet-to-cold transition rule: do not remain wet in wind/cold; change layers or substitute indoors.
  \end{itemize}
\item Hard-stop environment examples (non-negotiable):
  \begin{itemize}
  \item lightning/storm risk, ice glaze, whiteout/near-zero visibility, official air-quality warning, extreme windchill concern, heat index concern.
  \item posture: stop, substitute indoor modality, or reschedule to next valid window.
  \end{itemize}
\end{itemize}

\subsection*{7.14.6 Maintenance, Replacement, and Storage}
\phantomsection
\addcontentsline{toc}{subsection}{7.14.6 Maintenance, Replacement, and Storage}

\begin{itemize}
\item Drying posture:
  \begin{itemize}
  \item dry fully; avoid sealed storage while damp.
  \end{itemize}
\item Replacement cues:
  \begin{itemize}
  \item loss of water repellency that cannot be restored by care posture,
  \item seam failure or zipper failure,
  \item degraded insulation loft,
  \item glove wear causing cold exposure risk.
  \end{itemize}
\item Maintenance posture:
  \begin{itemize}
  \item wash per fabric requirements,
  \item re-treatment posture for water repellency as needed,
  \item seam repair tape as a field-reliability measure.
  \end{itemize}
\end{itemize}

\subsection*{7.14.7 Substitution Rules — Maintain Continuity Under Weather Constraints}
\phantomsection
\addcontentsline{toc}{subsection}{7.14.7 Substitution Rules — Maintain Continuity Under Weather Constraints}

\begin{itemize}
\item If conditions are unsafe:
  \begin{itemize}
  \item substitute indoor modality and document the environment deviation; do not force outdoor exposure.
  \end{itemize}
\item If apparel system is inadequate:
  \begin{itemize}
  \item downgrade exposure (duration/intensity) and substitute indoors if needed; do not “prove toughness.”
  \end{itemize}
\item Logging linkage (Stage 2.E posture, high-level):
  \begin{itemize}
  \item record environment notes when unsafe conditions drive substitution or re-slotting,
  \item record reason-code posture for the deviation,
  \item treat planned test attempts as compromised and reslot if conditions or apparel inadequacy would change performance tolerance.
  \end{itemize}
\end{itemize}

\subsection*{7.14.8 Phase Integration Map}
\phantomsection
\addcontentsline{toc}{subsection}{7.14.8 Phase Integration Map}

\begingroup
\scriptsize
\setlength{\tabcolsep}{2pt}
\renewcommand{\arraystretch}{1.12}

\begin{longtblr}{
  width=\linewidth,
  rowhead=1,
  colspec = {
    Q[l,wd=1.05in]
    Q[l,wd=1.55in]
    X[l,2.75]
    X[l,3.65]
  },
  hlines,
  vlines,
  cells = {hsep=2pt, vsep=6pt, halign=l, valign=t},
  row{odd} = {bg=white},
  row{even} = {bg=LightGray},
  row{1} = {bg=StageBlue!15, fg=black, font=\bfseries, halign=l, valign=m},
}
\Hdr{Phase} &
\Hdr{Required Tier (from Stage 6)} &
\Hdr{What Changes / Becomes Mandatory} &
\Hdr{Notes} \\
Phase 00 &
Tier 1 &
Cold: minimum layering core required; Wet: minimum shell posture required; Heat: minimum sun/heat posture required &
Phase 00 (End) requirement supports winter continuity. Do not introduce new apparel systems in protected windows. \\
Phase 01 &
Tier 1 &
Cold: maintain; Wet: maintain; Heat: maintain &
Phase 01 (End) emphasizes continuity; if winter execution is unstable, substitute indoors per Stage 7.08 rather than forcing. \\
Phase 02 &
Tier 1 &
Cold: maintain; Wet: maintain; Heat: maintain &
Phase 02 (End) maintain stable systems; avoid novelty near protected windows. \\
Phase 03 &
Tier 1 &
Cold: maintain; Wet: maintain; Heat: maintain &
Phase 03 (End) ruck capability begins to appear; avoid compounding novelty (new ruck interface + new cold system in same window). \\
Phase 04 &
Tier 1 &
Cold: maintain; Wet: maintain; Heat: maintain &
Phase 04 (End) aquatic capability begins; wet-gear handling is logistics-driven (Stage 7.10 / Stage 7.20), not "more apparel." \\
Phase 05 &
Tier 1 &
Cold: maintain; Wet: maintain; Heat: maintain &
Phase 05 (End) longer endurance increases consequence of exposure mismanagement; reslot rather than force. \\
Phase 06 &
Tier 2 &
Cold: expanded modularity required; Wet: packable storm kit required; Heat: improved heat management posture required &
Phase 06 (End) is first Tier 2 point. Traction aids and storm kit become operationally necessary for continuity in ice/wet conditions. \\
Phase 07 &
Tier 2 &
Cold: maintain expanded system; Wet: maintain storm kit; Heat: maintain heat posture &
Phase 07 (End) hybrid density increases consequence of cancellations; Tier 2 reduces reslots. \\
Phase 08 &
Tier 2 &
Cold: maintain; Wet: maintain; Heat: maintain &
Phase 08 (End) load tolerance increases fall risk under ice; traction aids remain critical. \\
Phase 09 &
Tier 2 &
Cold: maintain; Wet: maintain; Heat: maintain &
Phase 09 (End) high volume increases sweat risk; venting and drying posture become critical. \\
Phase 10 &
Tier 2 &
Cold: maintain; Wet: maintain; Heat: maintain &
Phase 10 (End) multi-domain gate weeks demand low variance; no novelty near protected windows. \\
Phase 11 &
Tier 2 &
Cold: maintain; Wet: maintain; Heat: maintain &
Phase 11 (End) recovery sensitivity increases consequence of exposure errors. \\
Phase 12 &
Tier 2 &
Cold: maintain; Wet: maintain; Heat: maintain &
Phase 12 (End) Stage 6 does not elevate 7.14 to Tier 3; overbuild remains optional and must not be implied as required. \\
Phase 13 &
Tier 2 &
Cold: maintain; Wet: maintain; Heat: maintain &
Phase 13 (End) consolidation demands repeatability; reslot if environment conditions compromise validity. \\
\end{longtblr}

\endgroup

7.14.8 Notes:

\begin{itemize}
\item Traction aids and packable storm kit are Tier 2-plus items; if any downstream artifact demands them earlier than Phase 06 (End), requires Stage 8 review.
\item If any phase card implies outdoor execution under hard-stop conditions is acceptable, requires Stage 8 review.
\end{itemize}

\subsection*{7.14.9 Risks, Contraindications, and Red Flags}
\phantomsection
\addcontentsline{toc}{subsection}{7.14.9 Risks, Contraindications, and Red Flags}

\begin{itemize}
\item Heat illness risk: overheating, dehydration, and sun exposure can cause unsafe sessions and invalid evidence; substitute indoors under heat index concern.
\item Cold injury risk: wet-to-cold failure, windchill exposure, and insufficient hand/head protection elevate risk; substitute indoors under extreme windchill concern.
\item Wet exposure risk: sustained wetness increases hypothermia risk and skin irritation; storm kit and drying posture reduce risk.
\item Slips/falls on ice: traction aids reduce risk but do not eliminate it; ice glaze remains a hard stop when unsafe.
\item Air quality risk: official air-quality warnings are hard-stop conditions; substitute indoors.
\item Visibility and traffic hazard posture: cross-reference Stage 7.19; do not attempt roadside exposure under near-zero visibility.
\item Validity contamination: tests attempted under unsafe or non-comparable conditions are inadmissible; reslot to a valid window.
\item Requires Stage 8 review triggers:
  \begin{itemize}
  \item any mismatch where a phase demands Tier 1/Tier 2 environmental apparel earlier than Stage 6 timing,
  \item any attempt to include anti-chafe consumables in Stage 7.14 tier tables (must cross-reference Stage 7.11 / Stage 7.15 / Stage 7.20),
  \item any attempt to treat hard-stop environment conditions as "train anyway" situations.
  \end{itemize}
\end{itemize}

\end{document}